\section{Interpretazione geometrica del determinante Jacobiano}

\begin{theorem}[significato geometrico Jacobiano]
  \label{th:significato_geometrico_jacobiano}
  Sia $L\colon \RR^m \to \RR^n$ una funzione lineare
  e sia $V$ l'immagine di $L$.
  Sia $M\colon V \to \RR^m$ una isometria.
  Sia $E\subset \RR^m$.
  Visto che le isometrie preservano le aree
  è naturale definire la $m$-area di $L(E)$ in $\RR^n$ 
  come la misura $m$-dimensionale di $M(L(E))$ in $\RR^m$
  Cioè per il teorema~\ref{th:significato_geometrico_determinante}
  \[
    \mathrm{m-area}(L(E)) = \abs{\det (M\circ L)}\cdot \abs{E}.
  \]
  Allora si ha 
  \[
    \mathrm{m-area}(L(E)) = \sqrt{\det (L^t L)}\cdot \abs{E}.
  \]
\end{theorem}
%
\begin{proof}
Estendiamo $M$ ad una isometria su tutto $\RR^n$, $\tilde M\colon \RR^n \to \RR^n$.
Dunque $\tilde M$ si rappresenta con una matrice ortogonale, cioè $\tilde M^t = \tilde M^{-1}$.
Sia $P\colon \RR^n \to \RR^m$ la proiezione ortogonale $P(x_1,\dots,x_n) = (x_1,\dots,x_m)$.
Dunque $M$ è la restrizione di $P\tilde M$ a $V$
e $M\circ L = P\tilde M L$.
Osserviamo che $P^t P\colon \RR^n\to \RR^n$ è l'applicazione 
lineare che azzera le ultime $n-m$ coordinate
ma visto che $M$ manda $V$ in $\RR^m$ si ha $P^t P \tilde ML = ML$. 
Dunque 
\[
 (\det(M\circ L))^2 
 = \det((P \tilde M L)^t (P \tilde M L)) 
 = \det(L^t \tilde M^t P^t P \tilde M L) 
 = \det(L^t \tilde M^t M L) = \det(L^t L)
\]
come volevamo dimostrare.
\end{proof}.

\begin{theorem}[formula di Cauchy-Binet]
  Siano $A$ una matrice $m\times n$ e $B$ una matrice $n\times m$.
  Allora 
  \[
    \det (AB) = \sum_{\vec k} \det A^{\vec k}\cdot \det B_{\vec k}
  \]
  dove 
  $\vec k = (k_1,\dots,k_m) \in \ENCLOSE{1,\dots,n}^m$ 
  è un elenco crescente di indici: $1\le k_1 < \dots < k_m\le n$
  e $A^{\vec k}$ e $B_{\vec k}$ sono le matrici $m\times m$ che 
  rappresentano i \emph{minori} di $A$ e $B$:
  \[
    (A^{\vec k})_{i,j} = A_{i,k_j}, \qquad
    (B_{\vec k})_{i,j} = B_{k_i, j}.
  \]
\end{theorem}
\begin{proof}
Denotiamo con $\Sigma_m$ il gruppo delle permutazioni 
di $\ENCLOSE{1,\dots,m}$ 
e con $(-1)^\sigma$ il segno della permutazione.
Si ha
\begin{align*}
  \det (AB) &= \sum_{\sigma \in \Sigma_m} (-1)^\sigma \prod_{i=1}^m (AB)_{i,\sigma(i)} \\
    &= \sum_\sigma (-1)^\sigma \prod_i \sum_{k=1}^n A_{i,k} B_{k,\sigma(i)} \\
    &= \sum_\sigma (-1)^\sigma \sum_{k_1=1}^n \dots \sum_{k_m=1}^n \prod_i A_{i,k} B_{k,\sigma(i)} \\
    &= \sum_\sigma \sum_{\vec k\in \ENCLOSE{1,\dots,n}^m} \prod_i A_{i,k} (-1)^\sigma \prod_i B_{k,\sigma(i)} \\    
    &= \sum_\sigma \sum_{\vec k\in \ENCLOSE{1,\dots,n}^m} \prod_i A_{i,k} \det B_{\vec k}
\end{align*}
Notiamo ora che se $\vec k = (k_1,\dots,k_m)$ non è iniettiva, cioè se $k_i = k_j$ per qualche 
$i\neq j$, la matrice $B_{\vec k}$ ha due righe uguali e quindi $\det B_{\vec k}=0$. 
Possiamo dunque restringere $\vec k$ alle funzioni iniettive.
Se $\vec k\colon\ENCLOSE{1,\dots,m}\to\ENCLOSE{1,\dots,n}$ è iniettiva c'è una 
unica permutazione $\phi$ tale che $i\mapsto k_{\phi(i)}$ è strettamente crescente.
Dunque ogni $\vec k$ iniettiva si ottiene prendendo tutte le $\vec k$ strettamente crescenti 
e componendole con tutte le permutazioni $\phi$ di $\ENCLOSE{1,\dots,m}$. 
Dunque
\begin{align*}
  \det(AB) 
    &= \sum_\sigma \sum_{1\le k_1<\dots<k_m\le n}\sum_{\phi\in \Sigma_m} 
    \prod_i A_{i,k_{\phi(i)}} \det B_{\vec k\circ \phi}\\
    &= \sum_\sigma \sum_{1\le k_1<\dots<k_m\le n}\sum_{\phi\in \Sigma_m} 
    \prod_i A_{i,k_{\phi(i)}} (-1)^\phi \det B_{\vec k}\\
    &= \sum_\sigma \sum_{1\le k_1<\dots<k_m\le n}\sum_{\phi\in \Sigma_m} 
    \det A^{\vec k} \det B_{\vec k}
\end{align*}
avendo usato il fatto che $\det B_{\vec k\circ \phi}$
si ottiene da $\det B_{\vec k}$ scambiando le righe 
secondo la permutazione $\phi$, ogni scambio di righe 
cambia segno al determinante e dunque si ha $\det B_{\vec k\circ \phi} = (-1)^\phi \det B_{\vec k}$.
\end{proof}


\section{integrale di superficie}

Dato $\phi\colon B\subset \RR^k \to A\subset \RR^n$
e data $f\colon A\subset \RR^n\to \RR$ 
possiamo definire
\[
  \int_\phi f(y) \, d\sigma(y) 
  \defeq \int_B f(\phi(x)) \sqrt{\det ((D\phi)^t D\phi)}\, dx.
\]

E' invariante per riparametrizzazioni.
Se $u\colon C\subset \RR^k \to B\subset \RR^k$ si ha 
\begin{align*}
  \int_{\phi\circ u} f(y) \, d\sigma(y) 
  &= \int_C f(\phi(u(x))) \sqrt{\det ((D(\phi\circ u))^t D(\phi\circ u))}\, dx \\
  &= \int_C f(\phi(u(x))) \sqrt{\det ((Du)^t (D\phi)^t D\phi Du)}\, dx \\
  &= \int_C f(\phi(u(x))) \sqrt{\det ((D\phi)^t D\phi)}\cdot \sqrt{\det ((Du)^t Du)}\, dx \\
  &= \int_B f(\phi(x)) \sqrt{\det ((D\phi)^t D\phi)}\, dx
\end{align*}
avendo fatto il cambio di variabile:
\[
  y = u(x), \qquad dy = \abs{\det Du} dx = \sqrt{\det ((Du)^t Du)}\, dx.
\]

Se $D\phi(x)$ ha rango massimo, i vettori $\phi_{x_1},\dots,\phi_{x_k}$
sono linearmente indipendenti e dunque 
generano uno spazio vettoriale $k$-dimensionale 
che si chiama spazio tangente alla superficie definita da $\phi$ nel punto $\phi(x)$.


\section{ipersuperifici}

Se $k=n-1$, $\phi\colon B\subset \RR^{n-1} \to A\subset \RR^n$ 
rappresenta una \emph{ipersuperficie} in $\RR^n$ o superficie di codimensione $1$.
Se $D\phi$ ha rango massimo in un punto $x$, 
i vettori $\phi_{x_1},\dots,\phi_{x_{n-1}}$ 
sono linearmente indipendenti e dunque
generano lo spazio tangente alla superficie definita da $\phi$ nel punto $\phi(x)$.
Per ogni $v\in \RR^n$ possiamo considerare il determinante $\det(v, D\phi)$ 
della matrice quadrata ottenuta aggiungendo la colonna $v$ alla matrice $D\phi$.
Chiamiamo $\xi_\phi$ il vettore che rappresenta l'applicazione lineare $v\mapsto \det(v, D\phi)$, 
cioè $\xi_\phi$ è definito univocamente dalla proprietà
\[
  \xi_\phi \cdot v = \det(v,D\phi).
\]
Chiaramente $\xi_\phi$ è ortogonale a tutte le colonne di $D\phi$ perché
$\xi_\phi \cdot \phi_{x_i} = \det(\phi_{x_i}, D\phi) = 0$,
dunque $\xi_\phi$ è un vettore ortogonale allo spazio tangente alla superficie.
Definiamo $\nu_\phi$ come il versore $\nu_\phi=\xi_\phi/\abs{\xi_\phi}$, 
che è uno dei due versori normali alla ipersuperficie.
La direzione di $\nu_\phi$ non cambia se si riparametrizza la superficie
e il suo verso non cambia se la riparametrizzazione ha jacobiano positivo.

Se $\phi$ non è iniettivo è possibile che la superficie si autointersechi 
o addirittura che si sovrapponga più volte su se stessa. 
In questo senso diremo che la superficie ha una \emph{molteplicità}.

Per quanto riguarda il modulo di $\xi_\phi$ si ha 
\begin{align*}
   \abs{\xi_\phi} &= \det (\nu_\phi, D\phi)
   = \sqrt{\det \enclose{(\nu_\phi, D\phi)^t (\nu_\phi,D\phi)}}  \\
   &= \sqrt{\det \begin{pmatrix}
     1 & 0 \\
     0 & (D\phi)^t D\phi
   \end{pmatrix}}
  = \sqrt{\det((D\phi)^t D\phi)}
\end{align*}
dunque 
\[
  \xi_\phi = \nu_\phi \sqrt{\det (D\phi)^t D\phi}
\]
e per ogni $v\in \RR^n$ si ha
\begin{equation}\label{eq:493444}
\det(v, D\phi)
= \xi_\phi \cdot v
= (\nu_\phi \cdot v) \sqrt{\det((D\phi)^t D\phi)}.
\end{equation}

Formalmente, se $\sigma$ rappresenta l'integrazione di superficie vista in precedenza, 
si ha
\[
  \xi_\phi(x) \, dx = \nu_\phi d\sigma.
\]


\subsection{flusso di un campo vettoriale attraverso una ipersuperficie}

Ad un campo vettoriale $\xi\colon A\subset \RR^n\to \RR^n$ possiamo associare la 
$(n-1)$-forma differenziale $\omega$ definita da
\[
  \omega_x[v_2,\dots,v_n] = \det(\xi(x), v_2,\dots,v_n).
\]
In tal modo, grazie a~\eqref{eq:493444} si ottiene che 
l'integrale di $\omega$ su una ipersuperficie rappresentata da $\phi\colon B\subset \RR^{n-1} \to \RR^n$
coincide con il flusso del campo vettoriale $\xi$ attraverso la ipersuperficie rappresentata da $\phi$: 
\begin{align*}
  \int_\phi \omega
  &= \int_B \det (\xi(x), D\phi(x))\, dx \\
  &= \int_B \xi(x) \cdot \nu_\phi(x) \sqrt{\det((D\phi)^t D\phi)}\, dx \\
  &= \int_\phi \xi \cdot \nu_\phi\, d\sigma.
\end{align*}
